\section{Toroidal field coils}
\label{sec:tf}
% Code: physics/compute_operating_params.m, wp_envelope.m,
% search/generate_combinations.m, scan_wp_designs.m, size_case_vault.m

\subsection{From configuration to operating parameters}
With $a=R_0/A$ and inboard build $w_b$, the inner-leg outer radius and the
outer-leg inner radius are
\begin{equation}
  R_1 = R_0-a-w_b,\qquad R_2=(R_0+a)\,\delta^{-1/N_\mathrm{TF}},
\end{equation}
where $\delta$ is the allowed ripple at the plasma edge
\todo{state the ripple approximation and its validity}. The current per
coil is $NI = 2\pi R_0 B_0/(\mu_0 N_\mathrm{TF})$ and the peak field on the
winding pack follows from the field at $R_1$ corrected for the finite
number of coils and for the winding-pack position
\todo{write the correction factor; NB the code uses
$2\pi R_1-\Delta r_\mathrm{ps}$ in the denominator: check whether it
should be $2\pi(R_1-\Delta r_\mathrm{ps})$}.

The coil is assumed to follow a bending-free (Princeton-D) shape between
$R_1$ and $R_2$, with $k=\tfrac12\ln(R_2/R_1)$
\todo{cite File, Mills \& Sheffield; Gralnick \& Tenney}. This gives in
closed form the centring-free hoop tension
$T=\tfrac12\,k\,N_\mathrm{TF}\,\mu_0 (NI)^2/(2\pi)$ and the inductance
\begin{equation}
L=\frac{\mu_0 (N_t k)^2 \sqrt{R_1R_2}}{2N_\mathrm{TF}}
\big[I_0(k)+2I_1(k)+I_2(k)\big].
\end{equation}
\note{This is where the global beam--shell model enters later: the
Princeton-D assumption gives membrane-only loads; the global model
removes it (section~\ref{sec:stellarator}).}

\subsection{Winding-pack layout and grading}
% - WP envelope: case width at R1, wedge (lateral) thickness range,
%   plasma-side case thickness, ground insulation, toroidal gaps
% - Turns per layer x number of layers -> NI/I_op within [Iop_min, Iop_max]
% - Field per grade assumed linear with the enclosed turns
% - 2 or 3 grades, grade boundaries at target fields
% - Turn reduction when the radial position shrinks the available width
%   (wedge geometry), extra layers added to recover the turns
\todo{text + Fig. 3: example WP cross-section with grades.}

\subsection{Case and inner vault}
The nose thickness is increased until both the vault and the innermost
jacket satisfy their Tresca allowables,
\begin{equation}
  \sigma_\mathrm{T,vault}=\sigma_z+\frac{2}{1-\beta^2}\,p_\mathrm{m}\,
  d_\mathrm{vault} \le \sigma_\mathrm{allow},\qquad
  \sigma_\mathrm{T,jacket}=\sigma_z+\sigma_\mathrm{rm}\,d_\mathrm{WP}
  \le \sigma_\mathrm{allow},
\end{equation}
with $\beta=R_k/R_1$, $\sigma_z=T/(A_\mathrm{jacket}+A_\mathrm{case})$ and
$d_\mathrm{vault}$, $d_\mathrm{WP}$ load-sharing factors from the
toroidal and radial stiffness of vault, case and winding pack.
\textbf{The wedged inner vault is what makes the tokamak TF tractable with
a 1D/2D model}: the centring force is reacted by the vault in toroidal
compression. This is precisely the assumption that fails for stellarator
coils.

\subsection{Design-space exploration}
% All combinations (wedge thickness x layouts) -> feasible table with
% sigma_T, I_op, J_eng, L, E, radial build, nose, WP h/w, turns, cable
% type, Cu/SC strands, hot-spot T, tau. Show a trade-off map (Fig. 4),
% e.g. radial build vs I_op coloured by J_eng, Pareto front.
\todo{Fig. 4: trade-off map for a reference machine.}
